\documentclass[12pt]{article}
\usepackage[utf8]{inputenc}
\usepackage[T1]{fontenc}
\usepackage{hyperref}
\title{Samatao Language Resources for K-12 Curriculum Development}
\author{Compiled by Claude Sonnet 5.5 (claude-sonnet-5-5) for Andr\'as Kornai}
\date{}
\begin{document}
\maketitle

\section{Introduction}
Samatao (ISO 639-3 ysd; Glottocode sama1295) is a Tibeto-Burman language of the Loloish (Ngwi) branch, spoken in and around Kunming, Yunnan, China. The community is officially counted within the Yi nationality. Estimates of speaker numbers vary widely. One linguistics paper reports that of nearly 3,000 registered Samatao people, fewer than 100, including semi-speakers, have active knowledge of the language. Verbix gives about 400 speakers. A Max Planck numerals page says there are no fluent speakers under 50. Urbanisation and Han migration are cited as causes of the decline. The ``Samtao'' profile on the Asia Harvest site is a different, Austro-Asiatic group and is not used here.

\section{Primary Educational Resources}
None found. I found no school programmes or textbooks.

\section{Linguistic and Reference Materials}
\begin{itemize}
\item Article in \emph{Linguistics of the Tibeto-Burman Area} 39(1) on Samatao endangerment (title not seen in the search results): \url{https://www.jbe-platform.com/content/journals/10.1075/ltba.39.1.06yu}
\item Samatao numeral data page (Max Planck, Leipzig): \url{https://lingweb.eva.mpg.de/channumerals/Samatao.htm}
\item Verbix endangered languages entry: \url{https://docs.verbix.com/EndangeredLanguages/Samatao}
\end{itemize}

\section{Religious and Community Resources}
\begin{itemize}
\item Joshua Project people-group page: \url{https://joshuaproject.net/people_groups/20613}
\end{itemize}

\section{Digital and Audio Resources}
\begin{itemize}
\item Global Recordings Network language page (I did not see whether recordings are actually available there): \url{https://globalrecordings.net/en/language/ysd}
\end{itemize}

\section{Curriculum Development Context}
There are few published resources and no known teaching materials. The language is critically endangered. Any curriculum would have to start with documentation and orthography work. That work should involve the community in Zijun and the surrounding Kunming area, and it should be coordinated with the authors of the academic literature above.
\end{document}